\chapter{Il modello a vincoli}
\label{ch:vincoli}

\epigraph{``Constraints are not the enemy of creativity. Constraints
are creativity's best friend.''}{Twyla Tharp,
\emph{The Creative Habit}}

\lettrine[lines=3]{C}{ostruire} un orario significa, dal
punto di vista formale, soddisfare un sistema di vincoli.
Alcuni di essi devono valere senza eccezione, altri esprimono
preferenze che il sistema cerca di rispettare il pi\`u
possibile pur sapendo che, in caso di conflitto, sar\`a
necessario un compromesso. Questo capitolo descrive con
ordine la tassonomia dei vincoli che $\pi$Tantum riconosce,
il modo in cui essi vengono espressi nell'interfaccia, e la
loro traduzione interna in vincoli per il solver CP-SAT.

La distinzione fra vincoli rigidi e morbidi non \`e un
dettaglio implementativo: \`e la chiave che permette al
sistema di restituire soluzioni utilizzabili. Trattare ogni
richiesta dell'utente come inderogabile produrrebbe sistemi
cronicamente infattibili, mentre trattare ogni vincolo come
puramente preferenziale violerebbe le regole contrattuali e
le esigenze didattiche di base. La gradazione in cinque
livelli che ora descriveremo \`e il risultato di un
compromesso meditato fra espressivit\`a per l'utente e
trattabilit\`a per il solver.

\section{La tassonomia dei cinque stati}

I vincoli di $\pi$Tantum sono organizzati in una tassonomia
uniforme di cinque stati, applicata sia alle matrici di
disponibilit\`a (cella per cella, per docenti, classi e
aule) sia alle preferenze di aula. I vincoli logici
disgiuntivi, espressi nel DSL del sistema, riconoscono
quattro di questi stati (l'ALLOWED non si applica a quel
contesto). La tabella~\ref{tab:cinque-stati} riassume i
cinque stati con il loro colore convenzionale, la loro
semantica e l'effetto che producono sul solver.

\begin{table}[h]
\centering
\sffamily\small
\begin{tabular}{l l p{8.5cm}}
\toprule
\textbf{Stato} & \textbf{Colore} & \textbf{Semantica} \\
\midrule
\textcolor{cAllow}{ALLOWED}  & verde chiaro & default neutro,
  applicabile soltanto alle griglie materia--aula e docente--aula;
  non contribuisce all'obiettivo \\
\textcolor{cSoft}{SOFT}      & giallo &
  penalit\`a positiva quando lo slot viene utilizzato; il valore
  \texttt{+penalty} viene sommato all'obiettivo che il solver
  minimizza \\
\textcolor{cHard}{HARD}      & rosso &
  vincolo inderogabile: il solver deve evitare lo slot oppure
  soddisfare il vincolo logico in modo assoluto \\
\textcolor{cPref}{PREFERRED} & blu &
  bonus negativo quando lo slot viene utilizzato o quando il
  vincolo logico \`e soddisfatto; il \texttt{-bonus} riduce
  l'obiettivo \\
\textcolor{cEnf}{ENFORCED}   & verde scuro &
  il solver deve forzare l'occorrenza dello slot o la
  soddisfazione del vincolo, con \emph{presence required}
  attiva \\
\bottomrule
\end{tabular}
\caption{I cinque stati della tassonomia di $\pi$Tantum, con
colori, semantica e contributo alla funzione obiettivo.}
\label{tab:cinque-stati}
\end{table}

Ognuno di questi stati ha un significato preciso che vale la
pena commentare. Lo stato ALLOWED, codificato in verde
chiaro, rappresenta la situazione di partenza: quella in cui
nessuna preferenza esplicita \`e stata espressa. \`E
applicabile soltanto a quelle griglie in cui ha senso
distinguere fra ``ammesso'' e ``vietato'', come la
compatibilit\`a fra una materia e un'aula. Per le matrici di
disponibilit\`a settimanale di docenti, classi e aule, lo
stato di partenza \`e implicitamente la disponibilit\`a, e
ALLOWED non viene mai utilizzato.

Lo stato SOFT, codificato in giallo, esprime una preferenza
negativa: l'utente preferirebbe non utilizzare quello slot,
ma se serve si pu\`o accettare. Quantitativamente, l'uso
dello slot aggiunge un costo positivo all'obiettivo, e il
solver cerca di evitarlo nei limiti del possibile. Il valore
del costo \`e configurabile, e l'utente lo regola attraverso
un campo numerico che appare nella cella stessa quando essa
viene marcata come gialla.

Lo stato HARD, codificato in rosso, esprime un divieto
assoluto. Il solver \emph{non pu\`o} utilizzare quello slot,
o equivalentemente \emph{non pu\`o} violare il vincolo
logico associato. Una soluzione che viola anche un solo
vincolo HARD non \`e una soluzione: viene scartata
immediatamente dal solver. Lo stato HARD si usa per
codificare le indisponibilit\`a contrattuali, le
incompatibilit\`a fisiche, e tutte le richieste che non
ammettono eccezioni.

Lo stato PREFERRED, codificato in blu, \`e simmetrico
rispetto al SOFT: esprime una preferenza positiva, una
configurazione che si vorrebbe vedere occorrere. L'uso
dello slot porta un bonus negativo all'obiettivo, e il
solver \`e quindi spinto a scegliere quella configurazione
quando ne ha possibilit\`a. \`E utile per esprimere
desideri come ``preferirei avere il laboratorio in due ore
consecutive'' o ``preferirei avere la palestra il
gioved\`i''.

Lo stato ENFORCED, codificato in verde scuro, esprime una
richiesta assoluta in senso opposto a HARD: il solver
\emph{deve} far accadere quella configurazione, non
soltanto deve permetterla. \`E lo strumento con cui si
codifica, per esempio, ``la riunione di consiglio si tiene
gioved\`i alla quinta ora'' (lo slot \`e
\emph{pre-allocato}). \`E un livello da usare con parsimonia,
perch\'e ENFORCED multipli e in conflitto fra loro rendono il
problema infeasibile, e il sistema non pu\`o fare altro che
segnalare l'impossibilit\`a.

\section{Le matrici di disponibilit\`a}

Le matrici di disponibilit\`a sono lo strumento principale
con cui l'utente esprime preferenze e indisponibilit\`a per
docenti, classi e aule. Ognuna di queste entit\`a possiede
una matrice di sei righe per sei colonne, che corrispondono
ai sei giorni della settimana scolastica e alle sei ore di
lezione tipiche di una giornata. Cliccando su una cella con
il mouse, l'utente la fa ciclare attraverso gli stati
disponibili nella sequenza \texttt{libero $\to$ giallo (SOFT)
$\to$ rosso (HARD) $\to$ blu (PREFERRED) $\to$ verde scuro
(ENFORCED) $\to$ libero}, in modo da poter raggiungere
qualunque stato con al massimo cinque click.

Quando una cella viene marcata come SOFT (gialla) o
PREFERRED (blu), accanto al colore appare un campo numerico
che permette di editare la penalit\`a o il bonus associati.
La logica dei segni \`e gestita automaticamente dal sistema:
i valori positivi sono ammessi soltanto sulle celle gialle e
i valori negativi soltanto su quelle blu, e il backend
applica un \emph{sign-clamp} al salvataggio per evitare
incoerenze. Trascinando il mouse mentre si tiene premuto il
tasto, si pu\`o applicare lo stesso stato a un blocco
contiguo di celle, comodit\`a che velocizza l'editing
quando un'intera giornata o un'intera fascia oraria deve
essere marcata in modo uniforme.

Una funzionalit\`a di sincronizzazione automatica governa il
rapporto fra il campo \emph{free day} del docente (il
giorno libero contrattuale) e la matrice. Quando il
\texttt{free\_day} \`e impostato, tutte e sei le ore di quel
giorno appaiono automaticamente HARD nella matrice, in modo
visivo, anche se l'utente non le ha marcate
esplicitamente. Viceversa, se l'utente edita la matrice
manualmente e tutte le ore di un giorno diventano HARD, il
backend deduce da questo che il \texttt{free\_day} \`e
quel giorno e lo aggiorna di conseguenza. La logica \`e
implementata nel modulo \texttt{webui/backend/db.py} e
garantisce che le due rappresentazioni restino sempre
coerenti.

La traduzione interna delle matrici nel modello CP-SAT
avviene nel modulo \texttt{optimization.py}, all'interno
della funzione \texttt{\_availability\_constraints}. Per
ogni entit\`a (docente, classe, aula), la matrice viene
analizzata e le sue celle si traducono in tre strutture
dati: un insieme di tuple \texttt{<entity>\_hard} che
rappresentano le \emph{forbidden assignments}, un dizionario
\texttt{<entity>\_soft} che mappa tuple a penalit\`a (con
segno gi\`a corretto), e un insieme di tuple
\texttt{<entity>\_enforced} che rappresentano le
\emph{presence required} nel solver. CP-SAT usa queste
strutture per costruire i vincoli durante la fase di
risoluzione.

\section{I vincoli logici disgiuntivi}

Accanto alle matrici di disponibilit\`a, il sistema offre
agli utenti uno strumento pi\`u espressivo per codificare
vincoli che non si lascino ridurre a singole celle: i
\emph{vincoli logici disgiuntivi}, spesso indicati con
l'acronimo DNF dall'inglese \emph{Disjunctive Normal Form}.
Si tratta di espressioni logiche che combinano slot e
predicati attraverso gli operatori AND, OR e NOT, e che si
applicano a una entit\`a (docente, classe, aula, indirizzo)
con un determinato livello di rigidit\`a (HARD, SOFT,
PREFERRED, ENFORCED).

\subsection{La grammatica delle espressioni}

La sintassi delle espressioni logiche di $\pi$Tantum \`e
stata pensata per essere compatta e leggibile. La forma BNF
semplificata \`e mostrata nel listato che segue, e descrive
gli elementi che si possono combinare.

\begin{lstlisting}[basicstyle=\ttfamily\small]
expr     := or_expr
or_expr  := and_expr ('OR' and_expr)*
and_expr := not_expr ('AND' not_expr)*
not_expr := ('NOT' | 'mai') not_expr | atom
atom     := slot | predicate | '(' expr ')'
slot     := <giorno><ora>
predicate:= <kind> ':' <name> ('@' <giorno> <ora>?)?
kind     := aula | materia | classe | gruppo | docente | prof
giorno   := lun|mar|mer|gio|ven|sab (anche lunedi|...|sabato,
            monday|...|saturday)
ora      := 1..6 (ordinale, 1=08:00) | 8..13 (assoluta)
\end{lstlisting}

La grammatica accetta anche le forme abbreviate \texttt{\&}
per AND, \texttt{|} per OR e \texttt{!} per NOT, oltre
all'alias italiano \texttt{mai} per NOT che produce
espressioni pi\`u naturali nei casi tipici. Per esempio,
\texttt{mai aula:LabFisica@mar} si legge in modo
trasparente come ``mai nel laboratorio di fisica il
marted\`i''.

\subsection{I predicate atoms}

Oltre agli slot temporali grezzi, le espressioni possono
contenere \emph{predicate atoms} che si riferiscono ad
entit\`a specifiche. La sintassi dei predicate atoms \`e
illustrata negli esempi seguenti.

\begin{lstlisting}[basicstyle=\ttfamily\small]
aula:LabFisica           -- usa l'aula LabFisica a qualunque slot
aula:LabFisica@mar       -- usa LabFisica il martedi a qualunque ora
aula:LabFisica@mar4      -- usa LabFisica martedi alla 4a ora (11:00)
materia:Religione@lun8   -- lezione di religione lunedi alle 8
classe:1A_Scientifico    -- riferimento alla classe 1A scientifico
gruppo:IRC               -- riferimento al gruppo IRC
\end{lstlisting}

I predicate atoms vengono parsati e persistiti correttamente
dal sistema. Il solver li tratta in modo coerente per le
restrizioni temporali; l'integrazione completa con la fase
di assegnazione delle aule \`e oggetto di sviluppo
incrementale, descritto nella roadmap del progetto.

\subsection{Esempi pratici}

Il modo migliore per familiarizzare con la sintassi dei
vincoli logici \`e leggere alcuni esempi reali di uso. I
seguenti esempi sono presi dalla pratica di alcune scuole
che usano $\pi$Tantum.

\begin{lstlisting}[basicstyle=\ttfamily\small]
% Per docenti
lun8 AND lun9
(lun8 AND lun9) OR (mar8 AND mar9)
gio11 OR gio12 OR gio13
NOT (mer3 AND mer4)
mai aula:LabFisica@mar
mai materia:Religione@lun8

% Per classi
mer4 AND mer5 AND mer6
mai aula:Palestra@gio
aula:LabInformatica@mar OR aula:LabInformatica@gio

% Per aule
gio4 AND gio5 AND gio6
mai materia:Religione
gruppo:IRC OR gruppo:Alternativa
\end{lstlisting}

Il primo esempio per i docenti, \texttt{lun8 AND lun9},
significa che il docente deve essere presente luned\`i alle
otto e alle nove (per esempio per una compresenza
obbligatoria). Il secondo esempio, con OR fra due
congiunzioni, ammette due alternative: o luned\`i alle
otto-nove, o marted\`i alle otto-nove. Il quarto, con NOT,
vieta la combinazione mercoled\`i terza-quarta ora. I tre
esempi con \texttt{mai} esprimono divieti su entit\`a
specifiche, leggibili in modo quasi colloquiale.

\subsection{Mapping CP-SAT}

Per ogni regola con DNF \texttt{[clause1, clause2, \ldots]},
il solver applica una traduzione che dipende dal kind del
vincolo. Una regola HARD impone che almeno una clausola sia
\emph{fully active}, ovvero che tutti i suoi atomi siano
veri. Una regola SOFT permette che nessuna clausola sia
attiva, ma in tal caso paga la penalit\`a positiva
\texttt{+soft\_penalty}. Una regola PREFERRED dota di un
bonus negativo \texttt{-soft\_penalty} la situazione in cui
almeno una clausola \`e attiva, riducendo cos\`i
l'obiettivo da minimizzare. Una regola ENFORCED si comporta
come una HARD ma con \emph{presence required}: almeno uno
degli slot della DNF deve coincidere con un'ora di lezione
attiva per l'entit\`a in questione, non basta che il vincolo
sia formalmente soddisfatto sui domini.

\section{I toggle HARD per classe}

Oltre ai vincoli espressi nelle matrici e ai vincoli logici,
il sistema mette a disposizione una serie di vincoli HARD
predefiniti, espressi come booleani nella tabella
\texttt{school\_classes}, che la maggior parte delle scuole
usa con valori standard. Sono sette: \texttt{hard\_entry\_at\_8}
(impone che la prima lezione cominci alle otto),
\texttt{hard\_exit\_after\_12} (impone che l'ultima lezione
non finisca prima delle dodici), \texttt{hard\_no\_holes}
(vieta i buchi nell'orario della classe),
\texttt{hard\_dual\_math} (impone che la matematica si
tenga in coppie di ore consecutive), \texttt{hard\_dual\_italian}
(idem per italiano), \texttt{hard\_motorie\_pairs} (idem per
educazione fisica), e \texttt{hard\_max\_6\_per\_day}
(impone il limite di sei ore al giorno).

A questi sette toggle si affianca il parametro SOFT
\texttt{soft\_minimize\_sixth\_weight}, che regola il peso
dell'obiettivo di evitare la sesta ora come ultima del
giorno. Il valore predefinito \`e tarato per esprimere una
preferenza moderata; valori pi\`u alti fanno s\`i che il
solver eviti la sesta ora con maggiore insistenza, valori
pi\`u bassi la accettano pi\`u liberamente.

\section{Il rilevamento dei conflitti}

Il tab Vincoli dell'interfaccia fornisce un bottone
denominato ``Cerca conflitti'' che, attivato, esegue una
diagnostica sul sistema di vincoli correntemente espresso e
segnala incoerenze rilevabili senza dover lanciare il
solver. Il bottone fa una chiamata HTTP a
\texttt{GET /api/monitor/conflicts}, e l'output viene mostrato
all'utente in una tabella con tre tipi di errore.

Il primo tipo, denominato \texttt{matrix\_hard\_enforced},
segnala le celle marcate sia come HARD sia come ENFORCED
contemporaneamente. Si tratta di una contraddizione
diretta, perch\'e una stessa cella non pu\`o essere allo
stesso tempo proibita e obbligatoria. Il secondo tipo,
denominato \texttt{enforced\_on\_free\_day}, segnala le
celle ENFORCED del docente che cadono nel suo giorno libero
contrattuale: anche questa \`e una contraddizione,
risolvibile o spostando l'ENFORCED o cambiando il
\texttt{free\_day}. Il terzo tipo, denominato
\texttt{logical\_unsatisfiable}, segnala i vincoli logici
HARD o ENFORCED la cui DNF non \`e soddisfacibile dato
l'insieme degli slot HARD del proprietario. \`E un caso che
richiede al coordinatore di rivedere il vincolo o di
rilassare le indisponibilit\`a.

\section{Tag delle aule}

Per esprimere caratteristiche delle aule che non si lasciano
ridurre al solo nome o tipo (``questa aula ha il
proiettore'', ``questa aula \`e adatta alla matematica
perch\'e ha la lavagna grande'', ``questa aula \`e una sala
riunioni''), il sistema offre i \emph{tag delle aule}: si
tratta di etichette libere, vere e proprie stringhe, che il
coordinatore associa a ciascuna aula dalla scheda Aule.
Ogni aula pu\`o avere quanti tag desidera, e ogni tag pu\`o
essere riutilizzato su pi\`u aule. La struttura non impone
una tassonomia gerarchica: i tag formano un semplice
insieme.

L'utilit\`a dei tag emerge soprattutto nella scrittura dei
vincoli. Invece di scrivere ``solo queste tre aule
specifiche'', enumerandole per nome (un elenco che andrebbe
aggiornato ogni volta che le aule cambiano), il coordinatore
pu\`o scrivere ``aule taggate matematica'', e il vincolo
seguir\`a automaticamente le etichettature correnti. Se
domani viene aggiunta una nuova aula taggata
\texttt{matematica}, il vincolo la includer\`a senza
ulteriori interventi.

Un esempio pratico chiarisce la logica. Supponiamo di
volere che la matematica si tenga soltanto in aule taggate
\texttt{matematica}. Una volta etichettate dalla scheda
Aule le aule che vogliamo includere, il vincolo nel DSL
generico si scrive cos\`i:

\begin{lstlisting}
forall l in lessons where l.subject == Matematica:
    "matematica" in l.classroom.tags
\end{lstlisting}

Il sistema offre anche una serie di tag automatici prodotti
dalla generazione mock. Per esempio, le aule home class
ricevono il tag dell'indirizzo della classe assegnata,
permettendo di scrivere vincoli del tipo ``le lezioni della
3A devono stare in aule taggate scientifico''. I dettagli
della generazione automatica si trovano nel modulo
\texttt{webui/backend/mock\_classrooms.py}.

\section{Le preferenze di giorno libero per i docenti}

I docenti hanno diritto, per contratto collettivo nazionale
del settore scolastico italiano, a un certo numero di giorni
liberi a settimana, tipicamente uno per il personale a
tempo pieno. $\pi$Tantum gestisce questa esigenza in modo
flessibile attraverso un meccanismo a tre slot di preferenza
ordinati. Il primo slot rappresenta il giorno libero
ideale del docente, per esempio il sabato; il secondo
rappresenta il giorno di ripiego se il primo non \`e
realizzabile, per esempio il mercoled\`i; il terzo
rappresenta l'ultima alternativa, per esempio il luned\`i.

Per ciascuna delle tre preferenze, il docente o il
coordinatore pu\`o scegliere se essa sia di tipo HARD
(assoluta, non negoziabile) o di tipo SOFT con un peso
configurabile. A questo si aggiunge il campo
\texttt{required\_free\_days\_count}, che indica quanti
giorni liberi servano in totale: il valore predefinito \`e
uno, conforme al contratto collettivo, ma scuole con
docenti a tempo parziale possono richiederne di pi\`u.

Un esempio illustra la flessibilit\`a del meccanismo. Si
consideri un docente con le preferenze [sabato HARD,
mercoled\`i SOFT con peso 50, luned\`i SOFT con peso 20] e
con \texttt{required\_free\_days\_count = 1}. Il sistema
cercher\`a sicuramente di mantenere libero il sabato,
perch\'e \`e HARD; se non riuscisse, in seconda istanza
cercher\`a di liberare il mercoled\`i (con un costo SOFT
di 50 se non ci riesce); in terza istanza cercher\`a di
liberare il luned\`i (con un costo di 20). La gradazione
permette al sistema di scegliere il compromesso pi\`u
ragionevole quando le preferenze multiple dei vari docenti
entrano in conflitto fra loro.

\section{Le preferenze di giorno libero per le classi}

Anche le classi possono avere giorni preferiti di non
attivit\`a, e la struttura del meccanismo \`e analoga a
quella dei docenti: tre slot ordinati con flag HARD o SOFT,
pi\`u un \texttt{required\_free\_days\_count}. La
motivazione \`e diversa: una classe pu\`o desiderare di
avere il sabato libero per attivit\`a di stage, oppure il
luned\`i per evitare il rientro tipicamente faticoso dopo il
fine settimana.

In aggiunta alle preferenze di giorno libero, le classi
possiedono il campo \texttt{max\_hours\_per\_day}, che
impone un tetto HARD al numero di ore in un giorno. Il
valore predefinito \`e cinque, ma alcuni istituti adottano
quattro per le classi del biennio (per ridurre il carico
giornaliero) o sei o sette per casi particolari. Riducendo
il tetto si ottengono giornate pi\`u leggere, ma se il
monte ore complessivo della classe \`e alto sar\`a
necessario estendere il calendario settimanale al sabato per
non sforare. Il sistema gestisce automaticamente questa
ridistribuzione, e segnala all'utente i casi in cui un
tetto troppo basso renderebbe il problema infeasibile.

\section{Compresenze, sostegno e potenziamento}

La scuola italiana ha tre famiglie di vincoli che il
sistema modella nativamente: le \textbf{compresenze} (il
laboratorio dove lavorano due docenti insieme), il
\textbf{sostegno} (il prof DVA che segue uno studente
specifico), e il \textbf{potenziamento} previsto dalla Legge
107 del 2015 (le ore dell'organico potenziato, senza
cattedra fissa). Sono entrati nel sistema a partire dal Task
C1 e sono enforced come constraint CP-SAT, non gestiti
post-hoc. Le violazioni sono restituite come errore al POST
con un messaggio specifico, prima che la run parta.

\subsection{Compresenze ``shared'' (lab di chimica)}

\textbf{Caso d'uso:} laboratorio di chimica del 2C, 4 ore
alla settimana, di cui 2 in compresenza con l'assistente di
laboratorio. La 2C ha le sue 4 ore di chimica, ma in 2 di
queste \`e prevista la presenza simultanea di due docenti.

\textbf{Configurazione dalla UI:}
\begin{enumerate}
\item Tab \texttt{/assignments}: la cattedra principale di
ProfChim (4 ore) \`e gi\`a presente.
\item Tab \texttt{/coteaching}: nuovo gruppo
\texttt{(2C, Chimica, n\_hours=2, kind=shared,
required=True)}.
\item Tab \texttt{/assignments}: aggiungi una nuova cattedra
ProfAss (2 ore di Chimica in 2C) e collegala al gruppo di
compresenza appena creato.
\end{enumerate}

\textbf{Cosa fa il solver:} introduce una variabile
\texttt{coday\_count[gid, d]} per ogni giorno, vincolata a
sommare a \texttt{n\_hours} sull'arco settimanale. Il
docente principale (quello con pi\`u ore, in questo caso
ProfChim) ha
\texttt{day\_count[principal, d] >= coday\_count[gid, d]}: le
sue ore di lezione coprono sempre quelle di compresenza. Il
co-docente (ProfAss) ha invece
\texttt{day\_count[codoc, d] == coday\_count[gid, d]}: le
sue 2 ore coincidono esattamente con le ore di compresenza.
A livello di slot orario, i due docenti vengono messi nella
stessa cella tramite una variabile booleana di slot
condiviso.

\subsection{Sostegno: il prof DVA che segue lo studente}

\textbf{Caso d'uso:} il prof Bianchi \`e il docente di
sostegno per uno studente DVA della 2A, e ha una cattedra
da 18 ore di sostegno. Il prof segue le lezioni della classe
(non ne aggiunge di proprie) e deve essere presente solo
quando 2A ha lezione con un altro docente.

\textbf{Configurazione:}
\begin{enumerate}
\item Tab \texttt{/teachers}: ProfBianchi esiste con la
materia ``sostegno'' nelle abilitazioni.
\item Tab \texttt{/assignments}: nuova cattedra
\texttt{(ProfBianchi, 2A, sostegno, ore=18,
is\_support=True)}.
\end{enumerate}

\textbf{Cosa fa il solver:} la cattedra di sostegno NON
contribuisce all'orario della 2A (le 18 ore non vanno a
sommarsi alle ore curricolari). Il vincolo \`e
\texttt{slot[ProfBianchi, 2A, sostegno, h] <= 2A\_busy[h]}:
il prof di sostegno \`e presente solo dove la 2A ha gi\`a
una lezione con un altro docente. La somma delle sue ore
settimanali \`e fissata a 18, distribuita in modo che il
profilo sia ragionevole.

\subsection{Potenziamento (Legge 107)}

\textbf{Caso d'uso:} la prof Verdi \`e nell'organico
potenziato e ha 18 ore settimanali di potenziamento, senza
una cattedra fissa. Le ore vengono usate per progetti,
recupero o supplenze.

\textbf{Configurazione:}
\begin{enumerate}
\item Tab \texttt{/assignments}: nuova cattedra
\texttt{(ProfVerdi, classe=NULL, materia=Potenziamento,
ore=18, is\_potenziamento=True)}. La cattedra viene
visualizzata in una sezione dedicata in fondo al tab,
\textbf{Potenziamento (Legge 107)}.
\end{enumerate}

\textbf{Cosa fa il solver:} introduce
\texttt{pot\_day\_count[ProfVerdi, d]} per ogni giorno, con
somma settimanale pari a 18. Il vincolo HARD
\texttt{pot\_day\_count[d] + ore\_di\_cattedra[d] <= 5}
garantisce che la prof non superi il tetto giornaliero. Cap
settimanale HARD: 30 ore (cinque giorni per cinque ore
ciascuno, escludendo il sabato).

\textbf{Visibilit\`a nel tab Assenze/Supplenze:} quando un
docente \`e assente, il sistema cerca un supplente fra i
disponibili. I docenti di potenziamento vengono mostrati per
primi, con un badge \textbf{POT} viola e un bordo dello
stesso colore: l'idea \`e che la loro disponibilit\`a sia
specificamente progettata per coprire questi buchi.

\section{Religione, alternativa e parallelismi intra-classe}

La classe 3B il marted\`i alla terza ora ha religione
cattolica con un docente, e contemporaneamente alternativa
con un altro docente, perch\'e gli studenti sono divisi.
Dal punto di vista dell'orario della 3B la classe \`e
``occupata'' in quello slot (gli studenti non sono
disponibili per altre lezioni), ma sono in atto due lezioni
parallele, non una.

\textbf{Configurazione (Task C2):}
\begin{enumerate}
\item Tab \texttt{/assignments}: due cattedre,
\texttt{(ProfRel, 3B, Religione, 1)} e
\texttt{(ProfAlt, 3B, Alternativa, 1)}.
\item Marca entrambe con lo stesso
\texttt{parallel\_group\_id} (id arbitrario; serve solo a
legarle insieme nel solver).
\end{enumerate}

\textbf{Cosa fa il solver:} le ore dei due docenti sono
legate slot-per-slot
(\texttt{slot[ProfRel, h] == slot[ProfAlt, h]}); a livello
di occupazione classe, le due lezioni condividono la stessa
chiave di busy, per cui la 3B risulta occupata UNA sola
volta in quello slot, anche se sono presenti due docenti.

\section{Gruppi inter-classe: seconda lingua, recupero}

Caso pi\`u complesso, introdotto dal Task C3: un gruppo che
attraversa pi\`u classi. Esempio: il gruppo di Spagnolo \`e
formato da 5 studenti della 2A e da 7 studenti della 2B,
si riunisce 3 ore alla settimana, ed \`e tenuto da
ProfSpa. Quando il gruppo si riunisce, le due classi-madre
2A e 2B sono ``occupate'' (i loro studenti sono nel gruppo,
non possono fare lezione in classe), anche se la lezione
non \`e formalmente una loro lezione.

\textbf{Configurazione:}
\begin{enumerate}
\item Tab \texttt{/students}: assicurarsi che gli studenti
del gruppo abbiano correttamente la classe-madre 2A o 2B.
\item Tab \texttt{/groups}: nuovo StudyGroup di nome
``Spagnolo 2A+2B'' con \texttt. Aggiungere i
12 studenti come membri e dichiarare le ore-materia con
\texttt{(Spagnolo, hours\_per\_week=3)}.
\item Tab \texttt{/assignments}: nuova cattedra,
selezionando il radio button \textbf{Target = gruppo
(StudyGroup)}, scegliere ``Spagnolo 2A+2B'' dal dropdown,
docente ``ProfSpa'', materia ``Spagnolo'', ore ``3''. La
cattedra appare in una sezione dedicata di colore ciano,
con badge \textbf{GRP}, separata sia dalle cattedre delle
classi che dalla sezione Potenziamento.
\end{enumerate}

\textbf{Cosa fa il solver:} per ogni ora in cui il gruppo
Spagnolo si riunisce, sia 2A che 2B vengono marcate come
occupate. L'invariante \texttt{somma di subj\_busy == pr}
del Phase B garantisce che, quando il gruppo \`e attivo in
uno slot, la classe-madre non possa avere lezioni regolari
nello stesso slot. Il vincolo per-day capacity in Phase A
\texttt{(cl\_day\_load + sum(group\_day\_count) <= 6)}
previene il caso degenere in cui curriculum e gruppo
sommerebbero pi\`u di sei ore in un solo giorno.

\subsection{Compresenza nel gruppo}

Se il gruppo Spagnolo prevede due docenti
(es. ProfSpa1 + ProfSpa2 in compresenza), si configura un
\texttt{CoteachGroup} con \texttt{group\_id} valorizzato
(XOR rispetto a \texttt{class\_id}) e si crea una seconda
cattedra con \texttt{coteach\_group\_id} sul gruppo. Il
modello CP-SAT applica la stessa logica principal/codoc che
gi\`a usa per le compresenze sulle classi.

\subsection{Sostegno nel gruppo}

Se uno studente DVA appartiene al gruppo Tedesco (3A+3B), il
prof di sostegno deve seguire lo studente \textbf{nel
gruppo} (non in classe-madre, perch\'e in quei tre slot lo
studente non \`e nella sua classe). Si configura una
cattedra
\texttt{(ProfSost, group\_id=tedesco, sostegno, ore=3,
is\_support=True)}. Il vincolo
\texttt{slot[ProfSost, h] <= group\_busy[h]} garantisce che
il prof sia presente solo nei tre slot in cui il gruppo si
riunisce, senza aggiungere ore di occupazione n\'e alla
classe-madre n\'e al gruppo.

\section{Lock nativi: blindare lezioni gi\`a sistemate}

Il pulsante a forma di lucchetto in \texttt{/assignments}
(e l'analogo nel monitor delle lezioni) marca una cattedra
come ``locked''. Quando il solver viene rilanciato per
ricalcolare l'orario, le lezioni di una cattedra lockata
sono trattate come constraint nativi: il sistema non pu\`o
spostarle da dove si trovano, deve risolvere intorno a
esse.

Vantaggio rispetto al vecchio meccanismo
``snapshot/restore'': il fail \`e immediato e parlante. Se
i lock sono incompatibili con altri vincoli (es. il giorno
libero del docente, il tetto giornaliero della classe, una
compresenza HARD), il sistema rifiuta la run con un errore
specifico, ad esempio:
\begin{quote}
\textit{Phase A INFEASIBLE: i lock attivi sono incompatibili
con i vincoli correnti. Rimuovi o adatta i lock (es.
sblocca eventi nel monitor) e ritenta. Lock attivi: 7
cattedra-giorno.}
\end{quote}

Senza lock nativi, lo stesso scenario sarebbe risultato in
una run apparentemente riuscita, con i lock silenziosamente
mossi a posteriori. I lock nativi sono propagati a tutte le
pipeline (monolitica, decomposizione temporale, spettrale,
curriculum, METIS, column generation) e a tutte le
metaeuristiche (LNS, SA, TS, ILS, ALNS, VNS, Lagrangian).

\section{Pre-flight checks: errori prima della run}

Quando l'utente lancia un'ottimizzazione, il sistema esegue
una serie di controlli sincroni \textbf{prima} di
schedulare la run. Le violazioni vengono restituite come
HTTP 400 con un elenco specifico, in italiano. Categorie
controllate:
\begin{itemize}
\item \textbf{Lock vs vincoli HARD}: lock che cadono nel
giorno libero del docente, oltre il tetto giornaliero della
classe, ecc.
\item \textbf{Compresenze}: numero di ore del principal
$\geq$ \texttt{n\_hours}; numero di ore del codoc esattamente
\texttt{n\_hours}.
\item \textbf{Sostegno}: deve avere una classe valida
(oppure un gruppo valido nel caso C3).
\item \textbf{Potenziamento}: \texttt{class\_id} deve essere
NULL; somma delle ore non deve superare 30.
\item \textbf{Gruppi (C3)}: XOR class\_id/group\_id; il
gruppo deve avere almeno uno studente; ogni studente deve
avere classe-madre; ore $>$ 0.
\end{itemize}

L'utente vede un toast di errore con la lista delle
anomalie. Pu\`o correggerle dai tab pertinenti
(\texttt{/assignments}, \texttt{/groups},
\texttt{/coteaching}) e rilanciare.

